% SPDX-License-Identifier: Apache-2.0
\section{Formalization}
\label{sec:formalization}

\subsection{The VWZ construction}

\claim{Describe VWZ: what it is, what it is for, and what security property it
is intended to provide. This must be written from the specification, not from
the Isabelle development, so that the formalization is seen to follow the
specification rather than define its own subject.}

The formalization is in Isabelle/HOL~\cite{isabelle}, on top of
\texttt{HOL-Number\_Theory}. The
session structure mirrors the dependency order of the mathematics:

\begin{center}
\begin{tabular}{@{}ll@{}}
\toprule
Session file & Contents \\
\midrule
\texttt{VWZ\_Fq3329} & the field $\mathbb{F}_q$ with $q = 3329$ \\
\texttt{VWZ\_Matrix} & matrices, determinant, elementary row operations \\
\texttt{VWZ\_Protocol} & protocol-level definitions \\
\texttt{VWZ\_Aux} & auxiliary lemma library \\
\texttt{VWZ\_ObligationsAligned} & the obligation ledger \\
\bottomrule
\end{tabular}
\end{center}

\subsection{Why tier the obligations}

A single ``proved / not proved'' bit is useless for a development that is
partially complete, and a bare \texttt{sorry} count is nearly as bad: it treats
an open goal that depends only on rewriting and an open goal whose subject has
not been defined as the same kind of debt. They are not. One is a proof waiting
to be written; the other is a definition waiting to be made, and no amount of
tactic engineering will close it.

We therefore tier obligations by \emph{what they depend on}:

\begin{center}
\begin{tabular}{@{}lrrl@{}}
\toprule
Tier & Lemmas & \texttt{sorry} & Character of the obligations \\
\midrule
T0 --- purely formal rewriting & 17 & 0 & closed \\
T1-A --- algebraic axioms & 16 & 0 & closed \\
T1-D --- definition-dependent & 5 & 2 & nearly closed \\
T2 --- linear algebra & 10 & 7 & substantially open \\
T3 --- security, probabilistic, parser & 8 & 8 & subjects undefined \\
\midrule
\textbf{Total} & \textbf{56} & \textbf{17} & \\
\bottomrule
\end{tabular}
\end{center}

The value of the split is that it makes the shape of the gap visible. ``17 open
goals'' sounds like a development that is 70\% done if one divides 56 into 39.
It is not, and the tier table says why: the closed tiers are the easy ones by
construction, and the open tiers are open for different reasons that require
different work.

\subsection{The closed tiers}

T0 covers purely formal rewriting: goals that follow from the definitions by
simplification and case analysis, with no mathematical content beyond
bookkeeping. The tier is not padding, and its lemmas are worth naming because
they show what ``closed'' buys. Three of them are unconditional equalities that
exist only to pin a definition down: \texttt{KeyGen\ k\ seed\ =\ KeyGen\ k\ seed}
and the same for \texttt{hash\_target} and \texttt{Sign}. They are trivial, and
they are recorded as obligations anyway, because a definition that cannot be
elaborated will not prove its own reflexivity, and that failure is worth
catching at the bottom of the stack rather than at the top.

The substantive T0 results are length and acceptance facts:

\begin{center}
\begin{tabular}{@{}ll@{}}
\toprule
Lemma & Statement (abbreviated) \\
\midrule
\texttt{T0\_B8} & serialised public key of the all-zero matrix has length $1 + 2nm^2$ \\
\texttt{T0\_C4\_size} & $n - w = k + 1$ \\
\texttt{T0\_C4\_ge3} & $2 \le k \Rightarrow 3 \le n - w$ \\
\texttt{T0\_E3a\_control\_accept} & the well-formed control instance is accepted \\
\texttt{T0\_E3a\_N1..N5} & five malformed instances are each rejected \\
\texttt{atMost\_Un\_atLeast\_nat} & $\{..c\} \cup \{c..\} = \mathrm{UNIV}$ \\
\texttt{ZaZb\_union} & $Z_a \cup Z_b = Z$ \\
\texttt{ZaZb\_inter} & $c \in Z \Rightarrow Z_a \cap Z_b = \{c\}$ \\
\bottomrule
\end{tabular}
\end{center}

The \texttt{T0\_E3a\_N1..N5} family is the interesting one. Four of the five are
negative: an empty public key, an empty signature, a mismatched $k$, and two
length off-by-ones. They are proved, not assumed, and together with the control
acceptance they establish that the length predicate is not the constant
\texttt{False} and not the constant \texttt{True}. That matters because a
predicate stuck at either constant would make every downstream parser lemma
vacuously true, and the development would look further along than it is.

T1-A covers the algebraic axioms of the field: commutativity and associativity
of addition and multiplication (\texttt{T1\_A1\_*}, \texttt{T1\_A2\_*}),
distributivity (\texttt{T1\_A3\_distrib}), additive and multiplicative identity
(\texttt{T1\_A4\_*}), inverses (\texttt{T1\_A5\_*}), the power laws
(\texttt{T1\_A6\_*}), the Fermat inverse $x^{-1} = x^{3327}$
(\texttt{T1\_A7\_inv\_fermat}), involution of inversion
(\texttt{T1\_A7\_inv\_inv}) and closure of the carrier
(\texttt{T1\_A8\_range}). These are instantiated from
\texttt{HOL-Number\_Theory.Residues} rather than reproved, and the obligation is
that the instantiation is the intended one: that \texttt{carrier},
$\oplus$ and $\otimes$ really are the field operations on $\mathbb{F}_{3329}$ and
not a structure that merely type-checks. \texttt{T1\_A7\_inv\_fermat} is the
load-bearing one, since the implementation computes inverses by exponentiation.
Two further T1-A lemmas are serialisation facts: \texttt{T1\_F2} fixes the
signature length at $1 + 4m$ and \texttt{T1\_F3\_u16le} fixes little-endian
16-bit encoding as $[x \bmod 256, \lfloor x/256 \rfloor \bmod 256]$.

\subsection{Linear algebra: what is done and what is not}

\texttt{VWZ\_Matrix} provides vectors and matrices as lists of lists, a
determinant defined by Laplace expansion along the first row, elementary row
operations, and rank-one and rank-two solvers.

The determinant is defined recursively:

\begin{lstlisting}
function det :: "int list list => int" where
  "det [] = 1"
| "det (r # rs) = sum_list (map (%j. (-1)^j * r!j *
                    det (map (del_col j) rs)) [0..<length r])"
\end{lstlisting}

The elementary row-addition matrix \texttt{row\_add\_matrix m i j f} is the
identity with $f$ placed at $(i,j)$. The obligation is that adding $f$ times row
$j$ to row $i$ does not change the determinant:

\begin{theorem}
For $0 < m$, $i < m$, $j < m$, $i \neq j$:
$\det(\texttt{row\_add\_matrix}\ m\ i\ j\ f) = 1$.
\end{theorem}

The case $i = 0$ is proved. The case $i > 0$ is \emph{not}, and
Section~\ref{sec:case-study-nat} explains exactly why: not because the
mathematics is in doubt, but because the proof requires an index shift that is
false at a boundary.

\subsection{Open obligations in T2 and T3}

The seven open T2 lemmas and the eight open T3 lemmas are listed below with
their statements. The T2 gaps are proof debt: the subjects are defined and the
statements are meaningful. The T3 gaps are mostly modelling debt: their subjects
are introduced as uninterpreted constants tagged \texttt{TODO(define)} with a
specification section reference, so there is not yet a formal statement about a
defined object to prove. Reporting both as a bare count without that distinction
would overstate the T3 work remaining.

\begin{center}
\begin{tabular}{@{}p{3.4cm}p{9.4cm}@{}}
\toprule
Lemma & Statement \\
\midrule
\multicolumn{2}{@{}l}{\textit{T2 --- proof debt}} \\
\texttt{T2\_B3} & $X_{2a}, X_{2b}, X_{3a}, X_{3b}$ are all invertible \\
\texttt{T2\_B5} & invertibility of $X_{2a}, X_{3a}$ gives $X_{2b}X_{2a}^{-1} = M_2$ and $X_{3b}X_{3a}^{-1} = M_3$ \\
\texttt{T2\_B7} & $\mathrm{rank}(\psi\text{-slice}) = 2$ \\
\texttt{T2\_D4} & $2 \le k \Rightarrow k - 1 \le \dim(\text{nullspace of } \texttt{rows3\_of}\ k)$ \\
\texttt{T2\_D5} & $i \in S_{\mathrm{set}} \Rightarrow \texttt{Eval}\ pk\ w_2\ w_3\, !\, i = \texttt{hash\_target}\ k\ msg\, !\, i$ \\
\texttt{T2\_D6} & $i \in Z_{\mathrm{set}} \Rightarrow P_{2a}(i) = 0 \lor P_{2b}(i) = 0$ \\
\texttt{T2\_D1\_correctness} & $\exists w_2 w_3.\ \texttt{Sign}\ sk\ msg = (w_2, w_3) \land \texttt{Eval}\ pk\ w_2\ w_3 = \texttt{hash\_target}\ k\ msg$ \\
\midrule
\multicolumn{2}{@{}l}{\textit{T3 --- modelling debt}} \\
\texttt{T3\_B1} & \texttt{distinct\_lam} returns $n$ distinct values in $\{1 \ldots 3328\}$ \\
\texttt{T3\_C1} & the Hamming weight of the target is $k + 1$ \\
\texttt{T3\_C3} & the sampler is unbiased \\
\texttt{T3\_D8\_empirical} & every message admits some signing pair \\
\texttt{T3\_E1\_parse\_pk\_total} & \texttt{parse\_pk} is total: it never diverges \\
\texttt{T3\_E2} & a signature of the wrong length is rejected \\
\texttt{T3\_E3b\_conjectural} & verification implies existence of a witness pair \\
\texttt{T3\_E4\_parse\_sig\_total} & \texttt{parse\_sig} is total \\
\bottomrule
\end{tabular}
\end{center}

One of these deserves a note. \texttt{T3\_E3b\_conjectural}
carries the word \emph{conjectural} in its name because it may not be true as
stated: it asserts that any accepted signature admits a witness pair, which is a
soundness-shaped claim about an uninterpreted \texttt{Verify}. We have kept it in
the ledger as a named conjecture rather than quietly weakening it, so that a
reader can see it is open for a stronger reason than a missing tactic. The T2
obligation that was open when this paper was first drafted, \texttt{T2\_B2}, has
since been closed; the obstruction it exposed and the way it was resolved are
Case~3 in Section~\ref{sec:case-study-nat}.

\subsection{The supporting library}

\texttt{VWZ\_Aux} contributes 49 lemmas and 1 theorem, none of which is open. It
is the largest single file in the development and contains the
index-arithmetic infrastructure that the determinant arguments need: head and
tail decomposition of sums, deletion of a column from a row, and the behaviour
of the determinant on rows with a single non-zero entry. It is described by
group rather than enumerated, because the individual statements are only
meaningful in terms of the matrix representation.

The groups are: row and column surgery (\texttt{del\_col} and its interaction
with \texttt{map}, \texttt{tl}, \texttt{replicate}); determinant behaviour on
structured rows (a row of zeros, a head row with one non-zero entry, a head row
that is a sum); split lemmas that decompose a row or a list at an index; and the
unit-vector lemmas that the elementary-matrix case analysis consumes.

Two named lemmas carry the T2 argument. \texttt{aux\_det\_id\_mat} states that
the determinant of the identity matrix is 1, proved by induction on the
dimension; without it the elementary-matrix obligation has no base case.
\texttt{aux\_det\_head\_one\_cons} is the head-collapse lemma: if a row begins
with 1 and all its remaining entries are 0, the determinant of the matrix headed
by that row equals the determinant of the minor obtained by deleting the first
column. This is the lemma that turns a Laplace expansion into an induction, and
Every case of the determinant argument runs through it. The theorem in this
file, \texttt{aux\_el\_det\_elementary\_i0}, handles the boundary case $i = 0$
that the induction cannot reach; it is a theorem rather than a lemma because the
ledger treats it as a top-level result for that case.

The remaining five lemmas exist because the $i > 0$ case needed them:
\texttt{aux\_t2\_nofire\_i} and \texttt{aux\_t2\_nofire\_j} cover the indices
at which the row operation does not fire at all, \texttt{aux\_t2\_head\_reduce}
expresses the determinant as the determinant of the head-deleted minor,
\texttt{aux\_t2\_row\_shift} is the shift lemma that is valid once $j = 0$ has
been excluded, and \texttt{aux\_t2\_minor0\_j0} / \texttt{aux\_t2\_minor0\_jpos}
give the two minor shapes separately. That split is the resolution described in
Section~\ref{sec:case-study-nat}.
